% !TEX root = ../AGFA-Light.tex

%% %%%%%%%%%%%%%%%%%%%%%%%%%%%%%%
%% -- ./content/AGFA-Section-2-Light.tex
%% -- Abschnitt 2 der Arbeit
%% -- Stand: 2026-09-24
%% %%%%%%%%%%%%%%%%%%%%%%%%%%%%%%
\section{Mein zweiter Abschnitt}\label{sec:zweiter-abschnitt}
\subsection{Test der Listen}\label{subsec:listen}
%%
Die Vorlage stellt mehrere Listenumgebungen bereit.
Der Name der Umgebung steht jeweils im ersten Eintrag.
%%
\begin{myenumerate}
\item
Aufzählung mit \verb|myenumerate|
\item
zweiter Eintrag
\end{myenumerate}
%%
\begin{myequivalent}
\item
Äquivalenzen mit \verb|myequivalent|
\item
zweite Aussage
\end{myequivalent}
%%
\begin{myitemize}
\item
Punkte mit \verb|myitemize|
\item
zweiter Punkt
\end{myitemize}
%%
\begin{mynumber}
\item
Nummeriert mit \verb|mynumber|
\item
zweiter Eintrag
\end{mynumber}
%%
Auch im laufenden Text ist eine Aufzählung möglich,
\begin{myinlinea}
\item
etwa mit \verb|myinlinea|
\item
wie hier,
\end{myinlinea}
wie dieses Beispiel zeigt.
%% --
\subsection{Test der mathematischen Umgebungen}\label{subsec:mumgebungen}
Seit über zweitausend Jahren gilt der Satz des Pythagoras als eines der schönsten Ergebnisse der Mathematik.
Wir beginnen mit einer kleinen Vorbereitung.
%
\begin{lemma}\label{lem:lemma1}
Ein rechtwinkliges Dreieck mit den Katheten $ a $ und $ b $ hat den Flächeninhalt $ \frac{ab}{2} $.
\end{lemma}
%
\begin{proof}
Zwei Exemplare des Dreiecks lassen sich entlang der Hypotenuse zu einem Rechteck mit den Seiten $ a $ und $ b $ zusammensetzen.
\end{proof}
%
\begin{theorem}[Pythagoras]\label{thm:theorem1}
In einem rechtwinkligen Dreieck mit den Katheten $ a $, $ b $ und der Hypotenuse $ c $ gilt stets
\begin{equation}\label{eq:pythagoras}
	a^{ 2 } + b^{ 2 } = c^{ 2 } .
\end{equation}
\end{theorem}
%
\begin{proof}
Man legt vier Exemplare des Dreiecks so in ein Quadrat der Seitenlänge $ a + b $, dass die Hypotenusen ein inneres Quadrat der Seitenlänge $ c $ begrenzen.
Vergleicht man die Flächeninhalte, so folgt mit \cref{lem:lemma1}
%
\[
	(a + b)^{ 2 } = c^{ 2 } + 4 \cdot \frac{ab}{2} ,
\]
%
also $ a^{ 2 } + 2ab + b^{ 2 } = c^{ 2 } + 2ab $ und damit \eqref{eq:pythagoras}.
\end{proof}
%%
\begin{corollary}\label{cor:folgerung}
Die Diagonale eines Quadrats mit der Seitenlänge $ 1 $ hat die Länge $ \sqrt{2} $.
\end{corollary}
%
\begin{proof}
Die Diagonale ist die Hypotenuse eines rechtwinkligen Dreiecks mit $ a = b = 1 $; nach \eqref{eq:pythagoras} gilt also $ c^{ 2 } = 2 $.
\end{proof}
%
\begin{proposition}[Umkehrung]\label{prop:prop}
Gilt für die Seiten $ a $, $ b $, $ c $ eines Dreiecks $ a^{ 2 } + b^{ 2 } = c^{ 2 } $, so ist das Dreieck rechtwinklig, und der rechte Winkel liegt der Seite $ c $ gegenüber.
\end{proposition}
%
\begin{remark}\label{rem:remark}
Der Satz des Pythagoras ist ein Spezialfall des Kosinussatzes
%
\[
	c^{ 2 } = a^{ 2 } + b^{ 2 } - 2ab \cos \gamma ,
\]
%
wobei $ \gamma $ der Winkel gegenüber der Seite $ c $ ist.
Für $ \gamma = \frac{\pi}{2} $ ergibt sich \eqref{eq:pythagoras}.
\end{remark}
%%
\begin{lem}\label{lem:lemma2}
Und noch ein Lemma, um zu testen, ob \texttt{lem} funktioniert.
\end{lem}
%%
Nun zur Eulerschen Zahl und zur imaginären Einheit:
%%
\begin{center}
\begin{tabular}{ll}
	\verb|$\eu$| 					& $ \eu $ \\
	\verb|$\iu$| 					& $ \iu $ \\
	\verb|$\eu^{\pi\iu} + 1 = 0$| 	& $ \eu^{ \pi \iu } + 1 = 0 $
\end{tabular}
\end{center}
%%
Die letzte Zeile ist die berühmte Eulersche Identität, die die fünf wichtigsten Konstanten der Mathematik verbindet:
%%
\begin{theorem}[Euler]\label{thm:eim}
Es gilt $ \eu^{ \pi \iu } + 1 = 0 $.
\end{theorem}
%%
Alles andere kann jeder selbst testen.
%
%%
\subsection{Querverweise}\label{subsec:referenzen}
\myquestion{Funktionieren alle Querverweise?}
%
Zunächst ein Verweis auf die Eulersche Identität mit \verb|\vref{thm:eim}|: \vref{thm:eim}.
Dann auf Gleichung~\eqref{eq:pythagoras} in \vref{thm:theorem1}.

Alternativ mit \verb|\cref|: \ldots \cref{thm:theorem1}, \cref{prop:prop}, \vref{eq:pythagoras} \ldots

Weitere Verweise: etwa auf \vref{lem:lemma1}, \vref{lem:lemma2}, \vref{cor:folgerung}, \vref{rem:remark} oder auf \vref{subsec:mumgebungen}.

Wie dies gemacht ist, sieht man in der \LaTeX-Datei.
%
\subsection{Sinnvolle Literatur zu \LaTeX}
%
Zum Nachschlagen aller Befehle und Möglichkeiten: \textcite{lshort-german} für den Einstieg, \textcite{latex-refsheet} als kompakte Übersicht und \textcite{l2tabu} für veraltete Befehle und Pakete, die man meiden sollte.
Die Links auf die Dokumente sind hinterlegt.

Weiter zu empfehlen sind \textcite{graetzer:2007} für den mathematischen Satz und \textcite{kohm:2014} für KOMA-Script, auf dem diese Vorlage aufbaut.
